\uuid{1rM9}
\exo7id{3902}
\titre{exo7 3902}
\auteur{quercia}
\organisation{exo7}
\datecreate{2010-03-11}
\isIndication{false}
\isCorrection{true}
\chapitre{Continuité, limite et étude de fonctions réelles}
\sousChapitre{Etude de fonctions}
\cpge{mpsi-mp2i.primitives-equations-differentielles.a, pcsi.techniques-calcul-primitives.a, ptsi.techniques-calcul-primitives.a, mpsi-mp2i.fonctions-valeurs-reelles-complexes.b, pcsi.techniques-calcul-fonctions.b, ptsi.techniques-calcul-fonctions.b, mpsi-mp2i.fonctions-valeurs-reelles-complexes.a, pcsi.techniques-calcul-fonctions.a, ptsi.techniques-calcul-fonctions.a, bcpst1.derivees-primitives-integrales.d, bcpst1.derivees-primitives-integrales.a, bcpst1.limites-continuite.c}
\module{Analyse}
\niveau{L1}
\difficulte{}

\contenu{
\texte{
Déterminer des primitives des fonctions suivantes :
}
\begin{enumerate}
    \item \question{$f(x) = \frac 1{\sqrt{x^2+x+1}}$.}
\reponse{$F(x) = \Argsh\left(\frac{2x+1}{\sqrt3}\right)$.}
    \item \question{$f(x) = \frac 1{x^2+x-1}$.}
\reponse{$F(x) = \frac 1{\sqrt5}\ln\left|\frac{\sqrt5 -2x-1}{\sqrt5 +2x+1}\right|$.}
\end{enumerate}
}
